\usepackage{graphicx,amsmath,amssymb,caption}
\newcommand{\setbaselineskip}[1]{\setlength{\baselineskip}{#1}}
% ===================== 参考文献 =====================
\usepackage[backend=biber,style=gb7714-2015]{biblatex}
\renewcommand{\bibfont}{\zihao{-5}\setbaselineskip{14bp}}
\setlength{\bibitemsep}{2bp plus 1pt}
\addbibresource{refs.bib}

% ===================== 配色 =====================
\usepackage{xcolor}
\definecolor{accent}{HTML}{1F4E8C}
\definecolor{accentlight}{HTML}{E8EEF6}
\definecolor{rowtint}{HTML}{F5F7FA}
\definecolor{inkmuted}{HTML}{55606E}
% 四段链路的固定配色，图与表一致
\definecolor{segbus}{HTML}{1F5FA8}
\definecolor{segrect}{HTML}{C8461E}
\definecolor{seghb}{HTML}{1E8A6E}
\definecolor{segpl}{HTML}{7A4FB0}

% ===================== 章节版式 =====================
\ctexset{
  chapter = {
    format     = \raggedright\heiti\zihao{3}\color{accent},
    nameformat = {},
    name       = {},
    number     = \arabic{chapter},
    aftername  = \hspace{0.8em},
    beforeskip = 0pt,
    afterskip  = 10bp,
    fixskip    = true,
  },
  section = {
    format    = \raggedright\heiti\zihao{-4}\color{accent},
    beforeskip= 10bp plus 2bp minus 1bp,
    afterskip = 4bp,
    aftername = \hspace{0.6em},
  },
  subsection = {
    format    = \raggedright\heiti\zihao{5},
    beforeskip= 6bp plus 1bp,
    afterskip = 2bp,
    aftername = \hspace{0.5em},
  },
}

% ===================== 版心 =====================
\usepackage{geometry}
\geometry{a4paper, top=23mm, bottom=20mm, left=22mm, right=22mm, headheight=21pt, headsep=6mm, footskip=8mm}
\setlength{\parskip}{0pt plus 1pt}

% ===================== 页眉页脚 =====================
\usepackage{fancyhdr}
\fancypagestyle{reportbody}{%
  \fancyhf{}%
  \fancyhead[L]{\small\color{inkmuted}托卡马克电源系统全链路仿真与校核}%
  \fancyhead[R]{\small\color{inkmuted}GKC3-Simulink 技术报告}%
  \fancyfoot[C]{\small\color{inkmuted}\thepage}%
  \renewcommand{\headrulewidth}{0.4pt}%
  \renewcommand{\headrule}{\hbox to\headwidth{\color{accent!35}\leaders\hrule height \headrulewidth\hfill}}%
  \renewcommand{\footrulewidth}{0pt}%
}

% ===================== 图 =====================
\graphicspath{{figures/}}
\DeclareGraphicsExtensions{.pdf,.png,.jpg}
\captionsetup{font=small, labelfont={bf,color=accent}, textfont=md, labelsep=quad, skip=3pt, justification=centering, singlelinecheck=true}
\captionsetup[table]{skip=4pt, position=top}
\setlength{\textfloatsep}{8pt plus 2pt minus 2pt}
\setlength{\floatsep}{7pt plus 2pt minus 2pt}
\setlength{\intextsep}{7pt plus 2pt minus 2pt}
\setcounter{topnumber}{3}
\setcounter{bottomnumber}{2}
\setcounter{totalnumber}{4}
\renewcommand{\topfraction}{0.92}
\renewcommand{\bottomfraction}{0.8}
\renewcommand{\textfraction}{0.08}
\renewcommand{\floatpagefraction}{0.8}
\usepackage{float}
\newcommand{\twofig}[4]{%
\begin{figure}[tbp]\centering
\includegraphics[width=0.8\linewidth,height=84mm,keepaspectratio]{#1}\\[0pt]{\small (a)}\\[2pt]
\includegraphics[width=0.8\linewidth,height=84mm,keepaspectratio]{#2}\\[0pt]{\small (b)}
\caption{#3}\label{#4}
\end{figure}}
% 模型视图分辨率足够，并排放置
\newcommand{\sidefig}[4]{%
\begin{figure}[tbp]\centering
\begin{minipage}[t]{0.49\linewidth}\centering\vspace{0pt}\includegraphics[width=\linewidth,height=58mm,keepaspectratio]{#1}\\[0pt]{\small (a)}\end{minipage}\hfill
\begin{minipage}[t]{0.49\linewidth}\centering\vspace{0pt}\includegraphics[width=\linewidth,height=58mm,keepaspectratio]{#2}\\[0pt]{\small (b)}\end{minipage}
\caption{#3}\label{#4}
\end{figure}}
\usepackage{placeins}
\makeatletter
\setlength{\@fptop}{0pt}
\setlength{\@fpsep}{8pt}
\setlength{\@fpbot}{0pt plus 1fil}
\makeatother

% ===================== 表 =====================
\usepackage{array,booktabs,tabularx,makecell,multirow,longtable}
\setlength{\aboverulesep}{0pt}
\setlength{\belowrulesep}{0pt}
\setlength{\heavyrulewidth}{0.9pt}
\setlength{\lightrulewidth}{0.45pt}
\arrayrulecolor{accent}
\newcommand{\tablefont}{\small\setlength{\extrarowheight}{0pt}\renewcommand{\arraystretch}{1.16}}
% 浅色表头与隔行底色
\newcommand{\headrow}{\rowcolor{accentlight}}
\newcommand{\zebra}{\rowcolors{1}{}{rowtint}}
\newcolumntype{L}{>{\raggedright\arraybackslash}X}
\newcolumntype{R}{>{\raggedleft\arraybackslash}X}
\newcolumntype{P}[1]{>{\raggedright\arraybackslash}p{#1}}
\newcolumntype{C}{>{\centering\arraybackslash}X}
\newenvironment{dtable}[2]{%
  \begin{table}[#1]\centering\tablefont\zebra\caption{#2}%
}{\end{table}}

% ===================== 正文小结框 =====================
\usepackage[most]{tcolorbox}
\newtcolorbox{keybox}{enhanced, breakable=false, colback=accentlight!55, colframe=accentlight!55,
  boxrule=0pt, arc=1.5pt, left=8pt, right=8pt, top=5pt, bottom=5pt, before skip=8pt, after skip=8pt,
  borderline west={2pt}{0pt}{accent}, fontupper=\small}

% ===================== 数学与单位 =====================
\usepackage{siunitx}
\sisetup{per-mode=symbol, detect-all=false, group-minimum-digits=5}
\providecommand{\dd}{\mathop{}\!\mathrm{d}}
\newcommand{\sub}[1]{_{\mathrm{#1}}}

% ===================== 绘图 =====================
\usepackage{tikz}
\usetikzlibrary{arrows.meta, positioning, fit, calc, backgrounds}
\usepackage[european resistors, american inductors]{circuitikz}

% ===================== 代码 =====================
\usepackage{listings}
\lstdefinestyle{repro}{%
  basicstyle=\ttfamily\footnotesize, columns=flexible, keepspaces=true, breaklines=true,
  backgroundcolor=\color{rowtint}, frame=leftline, rulecolor=\color{accent}, framerule=1.5pt,
  xleftmargin=1em, framexleftmargin=0.8em, aboveskip=6pt, belowskip=6pt,
  commentstyle=\color{inkmuted}, showstringspaces=false, upquote=true, extendedchars=false}
\lstnewenvironment{codeblock}[1][]{\lstset{style=repro,#1}}{}

% ===================== 超链接 =====================
\usepackage{hyperref}
\hypersetup{colorlinks=true, linkcolor=accent, citecolor=accent, urlcolor=accent,
  pdftitle={托卡马克电源系统全链路仿真与校核}, pdfauthor={GKC3-Simulink contributors}}
\urlstyle{tt}
\newcommand{\repourl}{https://github.com/KiloxGo/GKC3-Simulink}

% 正文行距
\AtBeginDocument{\setbaselineskip{17bp}}

% Allow long technical labels to wrap without changing report content.
\setlength{\emergencystretch}{3em}
